\documentclass[11pt, a4paper]{article}

% ── Codificación y Lenguaje ──────────────────────────────────────────────────
\usepackage[utf8]{inputenc}
\usepackage[T1]{fontenc}
\usepackage[spanish, es-tabla]{babel}

% ── Geometría ───────────────────────────────────────────────────────────────
\usepackage[top=2.5cm, bottom=2.5cm, left=2.5cm, right=2.5cm, headheight=24pt]{geometry}

% ── Matemáticas ─────────────────────────────────────────────────────────────
\usepackage{amsmath, amssymb, amsthm}
\usepackage{mathtools}

% ── Colores y Cajas ─────────────────────────────────────────────────────────
\usepackage{xcolor}
\usepackage[most]{tcolorbox}
\tcbuselibrary{skins, breakable, theorems}

% ── Tablas ──────────────────────────────────────────────────────────────────
\usepackage{booktabs}
\usepackage{array}
\usepackage{multirow}
\usepackage{longtable}

% ── Listas ──────────────────────────────────────────────────────────────────
\usepackage{enumitem}

% ── Encabezado y Pie de Página ──────────────────────────────────────────────
\usepackage{fancyhdr}

% ── Secciones ───────────────────────────────────────────────────────────────
\usepackage{titlesec}

% ── Gráficos y Diagramas ────────────────────────────────────────────────────
\usepackage{tikz}
\usetikzlibrary{shapes.geometric, arrows.meta, positioning, calc}
\usepackage{graphicx}

% ── Columnas múltiples (formulario) ─────────────────────────────────────────
\usepackage{multicol}

% ── Espaciado ───────────────────────────────────────────────────────────────
\usepackage{setspace}
\setlength{\parskip}{4pt}
\setlength{\parindent}{0pt}

% ── Paleta GOP ──────────────────────────────────────────────────────────────
\definecolor{gopblue}{RGB}{31, 78, 121}
\definecolor{gopcyan}{RGB}{70, 130, 180}
\definecolor{gopgreen}{RGB}{0, 100, 0}
\definecolor{gopgreenlight}{RGB}{230, 255, 230}
\definecolor{goporange}{RGB}{200, 100, 0}
\definecolor{goporangelight}{RGB}{255, 245, 210}
\definecolor{gopred}{RGB}{160, 0, 0}
\definecolor{gopredlight}{RGB}{255, 230, 230}
\definecolor{gopgray}{RGB}{90, 90, 90}
\definecolor{gopgraylight}{RGB}{245, 245, 240}
\definecolor{gopbluedef}{RGB}{0, 70, 140}
\definecolor{gopbluedeflight}{RGB}{220, 235, 255}

% ── Hipervínculos y TOC ─────────────────────────────────────────────────────
\usepackage[colorlinks=true, linkcolor=gopblue, urlcolor=gopblue]{hyperref}

% ── Entornos tcolorbox ───────────────────────────────────────────────────────
\newtcolorbox{definicion}[1]{
  enhanced, breakable,
  colback=gopbluedeflight, colframe=gopbluedef,
  fonttitle=\bfseries\small, title={Definicion: #1}, coltitle=white,
  attach boxed title to top left={yshift=-2mm, xshift=6pt},
  boxed title style={colback=gopbluedef, colframe=gopbluedef, sharp corners},
  sharp corners=south, top=4mm, before skip=6pt, after skip=6pt,
}
\newtcolorbox{formulas}[1]{
  enhanced, breakable,
  colback=gopgreenlight, colframe=gopgreen,
  fonttitle=\bfseries\small, title={Formulas: #1}, coltitle=white,
  attach boxed title to top left={yshift=-2mm, xshift=6pt},
  boxed title style={colback=gopgreen, colframe=gopgreen, sharp corners},
  sharp corners=south, top=4mm, before skip=6pt, after skip=6pt,
}
\newtcolorbox{tipprueba}[1]{
  enhanced, breakable,
  colback=goporangelight, colframe=goporange,
  fonttitle=\bfseries\small, title={TIP PRUEBA: #1}, coltitle=white,
  attach boxed title to top left={yshift=-2mm, xshift=6pt},
  boxed title style={colback=goporange, colframe=goporange, sharp corners},
  sharp corners=south, top=4mm, before skip=6pt, after skip=6pt,
}
\newtcolorbox{trampa}[1]{
  enhanced, breakable,
  colback=gopredlight, colframe=gopred,
  fonttitle=\bfseries\small, title={TRAMPA/ADVERTENCIA: #1}, coltitle=white,
  attach boxed title to top left={yshift=-2mm, xshift=6pt},
  boxed title style={colback=gopred, colframe=gopred, sharp corners},
  sharp corners=south, top=4mm, before skip=6pt, after skip=6pt,
}
\newtcolorbox{ejercicio}[1]{
  enhanced, breakable,
  colback=gopgraylight, colframe=gopgray,
  fonttitle=\bfseries\small\color{black}, title={Ejercicio: #1},
  attach boxed title to top left={yshift=-2mm, xshift=6pt},
  boxed title style={colback=gopgray!40, colframe=gopgray, sharp corners},
  sharp corners=south, top=4mm, before skip=6pt, after skip=6pt,
}
\newtcolorbox{resumentemario}{
  enhanced,
  colback=gopbluedeflight, colframe=gopbluedef,
  fonttitle=\bfseries, title={Resumen del Temario}, coltitle=white,
  attach boxed title to top left={yshift=-2mm, xshift=6pt},
  boxed title style={colback=gopbluedef, colframe=gopbluedef, sharp corners},
}
\newtcolorbox{estructuraprueba}{
  enhanced,
  colback=goporangelight, colframe=goporange,
  fonttitle=\bfseries, title={Estructura de la Prueba}, coltitle=white,
  attach boxed title to top left={yshift=-2mm, xshift=6pt},
  boxed title style={colback=goporange, colframe=goporange, sharp corners},
}
\newtcolorbox{memorizar}{
  enhanced, breakable,
  colback=gopredlight, colframe=gopred,
  fonttitle=\bfseries, title={Conceptos Clave --- Memorizar esto}, coltitle=white,
  attach boxed title to top left={yshift=-2mm, xshift=6pt},
  boxed title style={colback=gopred, colframe=gopred, sharp corners},
  sharp corners=south,
}
\newtcolorbox{paradesarrollo}{
  enhanced, breakable,
  colback=gopgreenlight, colframe=gopgreen,
  fonttitle=\bfseries, title={Para la Etapa Desarrollo (Apuntes Libres Permitidos)}, coltitle=white,
  attach boxed title to top left={yshift=-2mm, xshift=6pt},
  boxed title style={colback=gopgreen, colframe=gopgreen, sharp corners},
  sharp corners=south,
}

% ── Encabezado y pie ─────────────────────────────────────────────────────────
\pagestyle{fancy}
\fancyhf{}
\fancyhead[L]{\small\textbf{GOP ICS3213} --- Estudio \texttt{I2} \texttt{2026}}
\fancyhead[C]{}
\fancyhead[R]{\small\nouppercase{\rightmark}}
\fancyfoot[L]{\footnotesize Generado a partir de clases, ayudantías 6--8, I2 2025 y casos 2026}
\fancyfoot[R]{\small\thepage}
\renewcommand{\headrulewidth}{0.4pt}
\renewcommand{\footrulewidth}{0.4pt}

% ── Estilo de secciones ──────────────────────────────────────────────────────
\titleformat{\section}
  {\color{gopblue}\Large\bfseries}{\color{gopblue}\thesection.}{0.5em}{}
  [\color{gopblue}\titlerule]
\titleformat{\subsection}
  {\color{gopblue}\normalsize\bfseries}{\color{gopblue}\thesubsection.}{0.5em}{}
\titleformat{\subsubsection}
  {\color{gopblue}\small\bfseries}{\color{gopblue}\thesubsubsection.}{0.5em}{}

\begin{document}

\begin{titlepage}
  \centering
  \vspace*{2cm}
  {\Huge\bfseries\color{gopcyan} Estudio \texttt{I2} - 04 Proyectos\par}
  \vspace{0.4cm}
  {\LARGE\bfseries Gestión de Operaciones\par}
  \vspace{0.3cm}
  {\large ICS3213 --- PUC Chile --- \texttt{2026}\par}
  \vspace{1.2cm}
  \rule{\textwidth}{0.4pt}
  \vspace{0.4cm}

  {\normalsize \texttt{Miércoles 3 de Junio de 2026, 17:30 hrs (presencial)}\par}
  \vspace{0.2cm}
  {\small Prof.\ Alejandro Mac Cawley --- Rodrigo Carrasco\par}
  \vspace{0.2cm}
  {\small\textbf{Autor:} \href{https://cl.linkedin.com/in/fabian-ortega-llanten}{Fabián Ignacio Ortega Llantén}\par}
  \vspace{0.4cm}
  \rule{\textwidth}{0.4pt}
  \vspace{1cm}

  \begin{resumentemario}
    \begin{itemize}[leftmargin=1.2em, itemsep=2pt]
      \item Planificación Agregada de la Producción
      \item Planificación de Corto Plazo y MRP
      \item Administración de Proyectos (CPM / PERT)
      \item Variabilidad en las Operaciones (Teoría de Colas)
      \item Bodegas, Centros de Distribución y Decisiones Estratégicas
      \item \textbf{La Meta:} Capítulos 16 al 26 (inclusive)
      \item \textbf{Caso 3:} Barilla SpA \quad y \quad \textbf{Caso 4:} University Health System
      \item \textbf{Juego de la Cerveza} (efecto látigo / bullwhip)
      \item Cápsulas: Ejercicios de Planificación de la Producción y planilla MRP
    \end{itemize}
  \end{resumentemario}

  \vspace{0.5cm}

  \begin{estructuraprueba}
    \begin{itemize}[leftmargin=1.2em, itemsep=3pt]
      \item \textbf{Etapa Digital (60 min, SIN apuntes):}
        \begin{itemize}[leftmargin=1.2em, itemsep=1pt]
          \item 10 preguntas selección múltiple (materia)
          \item 10 preguntas V/F de La Meta (si es Falso, argumentar por qué)
          \item 1 pregunta de desarrollo
        \end{itemize}
      \item \textbf{Descanso de 5 minutos}
      \item \textbf{Etapa Desarrollo (60 min, APUNTES FÍSICOS + formulario):}
        \begin{itemize}[leftmargin=1.2em, itemsep=1pt]
          \item Ejercicios de cálculo (Planif. Agregada, MRP, PERT, Colas)
          \item Apuntes solo en papel; NO tablets ni celulares.
          \item Escanear y subir a Canvas; dejar prueba física en el banco.
        \end{itemize}
    \end{itemize}
  \end{estructuraprueba}

  \vfill
\end{titlepage}
\newpage
\section{Módulo 3: Administración de Proyectos (CPM / PERT)}

\subsection{Conceptos y triángulo de hierro}

\begin{definicion}{Proyecto y sus objetivos}
Un \textbf{proyecto} es un esfuerzo temporal con inicio y fin para crear un resultado único. Sus tres objetivos forman el \textbf{triángulo de hierro}: \textbf{Alcance, Tiempo y Costo}. Están en conflicto permanente: mejorar uno típicamente empeora otro (\textit{trade-off}). El ciclo de vida: Concepto $\to$ Factibilidad $\to$ Planificación $\to$ Ejecución $\to$ Término.
\end{definicion}

\begin{definicion}{Definiciones clave de redes}
\begin{itemize}[leftmargin=1.4em, itemsep=1pt]
  \item \textbf{Pseudo-actividad (dummy):} actividad de duración cero, no consume recursos; solo indica precedencias.
  \item \textbf{Ruta crítica:} secuencia de actividades que forman la \textbf{cadena más larga} en tiempo. SIEMPRE existe al menos una; puede haber varias.
  \item \textbf{Holgura (slack):} margen de retraso de una actividad \textbf{sin} afectar el fin del proyecto. Las actividades de la ruta crítica tienen holgura \textbf{cero}.
\end{itemize}
\end{definicion}

\subsection{CPM: el método de la ruta crítica}

\begin{formulas}{Cálculo de tiempos CPM (forward y backward pass)}
\begin{align*}
  EF &= ES + t && \text{(término más temprano)}\\
  ES_j &= \max\{EF_{ij}\} && \text{(inicio más temprano = máx de predecesores)}\\
  LS &= LF - t && \text{(inicio más tardío)}\\
  LF_k &= \min\{LS_{kj}\} && \text{(término más tardío = mín de sucesores)}\\
  H &= LF - EF = LS - ES && \text{(holgura)}
\end{align*}
\textbf{Forward pass} (izq.\ a der., $ES/EF$): la duración del proyecto $T$ es el mayor $EF$. \textbf{Backward pass} (der.\ a izq., $LS/LF$): en el nodo final $LF=T$. \textbf{Ruta crítica:} camino continuo de actividades con $H=0$.
\end{formulas}

\subsection{PERT: incorporando incertidumbre}

\begin{formulas}{PERT con 3 estimaciones (distribución Beta)}
Cada actividad tiene tres estimaciones: $a$ (optimista), $m$ (más probable), $b$ (pesimista). Bajo distribución Beta:
\[
  t_e = \mu = \frac{a + 4m + b}{6}, \qquad
  \sigma = \frac{b - a}{6}, \qquad
  \sigma^2 = \left(\frac{b-a}{6}\right)^2
\]
\end{formulas}

\begin{formulas}{Duración del proyecto y probabilidades}
Por el \textbf{Teorema Central del Límite}, la duración total $T$ (suma de actividades de la ruta crítica) se distribuye \textbf{Normal}:
\[
  T \sim \mathcal{N}\!\left(\textstyle\sum_{i\in RC} t_{e,i},\ \sum_{i\in RC} \sigma_i^2\right)
\]
Para calcular probabilidades se estandariza:
\[
  z_0 = \frac{t_0 - T}{\sigma_T}, \qquad \sigma_T = \sqrt{\textstyle\sum_{i\in RC}\sigma_i^2}
\]
\textbf{Intervalo de confianza} al $(1-\alpha)$: $\;T \pm z_{\alpha/2}\cdot\sigma_T$.
\end{formulas}

\begin{trampa}{La varianza SOLO suma sobre la ruta crítica}
$\sigma_T^2$ se calcula sumando las varianzas \textbf{únicamente} de las actividades de la ruta crítica, no de todas. Si una pregunta pide la probabilidad de que \textbf{otra} ruta (no crítica) se vuelva crítica, debes calcular el tiempo y varianza \textbf{de esa ruta} y, asumiendo independencia, $P(\text{esa ruta} > T_{RC})$. Para la probabilidad de que una sea la crítica entre varias, se multiplican las probabilidades (independencia).
\end{trampa}

\subsection{Modelos tiempo-costo (Crashing)}

\begin{formulas}{Crashing: acortar la ruta crítica al mínimo costo}
Notación: $NT$ (tiempo normal), $CT$ (tiempo acelerado/crash), $NC$ (costo normal), $CC$ (costo crash). El costo de acelerar por unidad de tiempo (pendiente):
\[
  \text{Slope} = \frac{CC - NC}{NT - CT}
\]
\textbf{Procedimiento:} (1) preparar diagrama CPM, (2) calcular slope de cada actividad, (3) hallar ruta crítica, (4) \textbf{acortar la actividad crítica de menor slope}, un paso a la vez. ¡Ojo! al acortar, la ruta crítica puede cambiar.
\end{formulas}

\begin{tipprueba}{Bono y penalización (pregunta clásica de PERT)}
\textbf{Identificación:} te ofrecen un bono \$$X$ si terminas en $\le t_1$ días y/o penalización \$$Y$ si terminas en $\ge t_2$. \textbf{Qué aplicar:} calcula el \textbf{valor esperado} de cada incentivo: $VE_{\text{bono}} = P(T\le t_1)\cdot X$ y $VE_{\text{pen}} = P(T\ge t_2)\cdot Y$. Aceptas el contrato si $VE_{\text{bono}} > VE_{\text{pen}}$. Para el punto de \textbf{indiferencia}, igualas $VE_{\text{bono}} = P(T\ge t^*) \cdot Y$ y despejas $t^*$ vía la tabla normal.
\end{tipprueba}

\begin{tipprueba}{Crashing con bono: comparar escenarios por VE neto}
Cuando el crashing se combina con un bono por terminar antes, no basta con ``minimizar el costo de crashing''. El procedimiento correcto es:
\begin{enumerate}[leftmargin=1.5em, itemsep=2pt]
  \item Para cada nivel de reducción posible (0, 1, 2, \ldots semanas), identificar la combinación más barata de actividades a acortar en la ruta crítica.
  \item Para cada escenario, calcular $VE_{\text{neto}} = \text{bono} \cdot P(T < T_{\text{bono}}) - \text{costo de aceleración}$.
  \item Elegir el escenario con mayor $VE_{\text{neto}}$. ¡Puede ser 0 semanas (no acelerar)!
  \item \textbf{Verificar rutas alternativas:} tras cada reducción, calcular la duración de las rutas no críticas. Si alguna supera la nueva RC, se convierte en la nueva ruta crítica y debe incluirse en el siguiente paso.
\end{enumerate}
\textit{Error frecuente:} resolver directamente ``¿cómo llego de 20 a 18 semanas al mínimo costo?'' sin evaluar si acortar 1 semana en vez de 2 da mayor VE neto.
\end{tipprueba}


\clearpage
\subsection[Proyectos: PERT y Crashing (Ay. 8)]{Ejercicio de Aplicación: PERT Completo, Crashing y Evaluación de Contrato (Ayudantía 8)}

\begin{ejercicio}{PERT estocástico, IC y crashing de costo mínimo (Ayudantía 8)}
\begin{tipprueba}{¿Cuándo usar este enfoque?}
\textbf{Identificación:} te dan actividades, dependencias, y tiempos estimulados (optimista, más probable, pesimista). Piden calcular $t_e$ y $\sigma^2$ por actividad, dibujar diagrama, encontrar ruta crítica, calcular probabilidades normales e intervalos de confianza, y determinar el crashing óptimo.
\textbf{Qué aplicar:} Fórmulas Beta de PERT, método de ruta crítica (holguras = 0), Teorema Central del Límite para el proyecto ($T_e = \sum_{RC} t_{e,i}$, $\sigma_T^2 = \sum_{RC} \sigma_i^2$), normalización $Z$ y análisis de costo marginal de crashing.
\end{tipprueba}

\textbf{Enunciado:}
Se dispone de la relación de dependencia y tiempos (en semanas) de las etapas de un proyecto:
\begin{center}
\small
\begin{tabular}{@{}ccccc@{}}
  \toprule
  \textbf{Etapa} & \textbf{Predecesor} & \textbf{Optimista ($a$)} & \textbf{Probable ($m$)} & \textbf{Pesimista ($b$)} \\
  \midrule
  A & - & 2 & 3 & 4 \\
  B & A & 2 & 4 & 6 \\
  C & A & 5 & 6 & 13 \\
  D & B, C & 3 & 6 & 9 \\
  E & B & 2 & 5 & 8 \\
  F & D, E & 2 & 4 & 6 \\
  \bottomrule
\end{tabular}
\end{center}
1. Encuentre los tiempos esperados y varianzas de cada actividad.
2. Obtenga los tiempos ES, EF, LS, LF y holgura para cada etapa. Indique la ruta crítica y duración esperada.
3. Calcule la probabilidad de terminar en menos de 22 semanas, y el intervalo de confianza al 95\%.
4. Si el cliente ofrece un bono de \$8.000 al finalizar el proyecto si este se termina en menos de 18 semanas, ¿qué actividades acortaría usando la tabla de costos de crashing?
\begin{center}
\small
\begin{tabular}{@{}cccc@{}}
  \toprule
  \textbf{Actividad} & \textbf{Reducción Máx. (semanas)} & \textbf{Costo Normal (\$)} & \textbf{Costo Acelerado (\$)} \\
  \midrule
  A & 1 & 10.000 & 13.000 \\
  B & 1 & 6.000 & 9.000 \\
  C & 2 & 4.000 & 7.000 \\
  D & 2 & 13.000 & 18.000 \\
  E & 2 & 9.000 & 13.000 \\
  F & 1 & 7.000 & 8.000 \\
  \bottomrule
\end{tabular}
\end{center}

\textbf{Solución:}

\textbf{1) Tiempos esperados y varianzas:}
Usando $t_e = \frac{a+4m+b}{6}$ y $\sigma^2 = \left(\frac{b-a}{6}\right)^2$:
\begin{itemize}
  \item \textbf{A:} $t_e = 3$ semanas; $\sigma^2 = 1/9 \approx 0,111$.
  \item \textbf{B:} $t_e = 4$ semanas; $\sigma^2 = 4/9 \approx 0,444$.
  \item \textbf{C:} $t_e = 7$ semanas; $\sigma^2 = 16/9 \approx 1,778$.
  \item \textbf{D:} $t_e = 6$ semanas; $\sigma^2 = 1,0$.
  \item \textbf{E:} $t_e = 5$ semanas; $\sigma^2 = 1,0$.
  \item \textbf{F:} $t_e = 4$ semanas; $\sigma^2 = 4/9 \approx 0,444$.
\end{itemize}

\textbf{2) Tiempos de calendario (Forward and Backward Pass):}

Para calcular los tiempos de calendario de cada actividad, aplicamos recursivamente las siguientes reglas:
\begin{itemize}
  \item \textbf{Forward Pass (Tiempos Tempranos):} Recorremos el proyecto desde el inicio hasta el final para determinar el inicio más temprano ($ES$) y el término más temprano ($EF$).
    \begin{itemize}
      \item Para actividades sin predecesores (iniciales): $ES = 0$.
      \item Para cualquier otra actividad $i$: $ES_i = \max_{j} (EF_j)$ de todos sus predecesores inmediatos $j$.
      \item Término temprano: $EF_i = ES_i + t_e(i)$.
      \item La duración esperada total del proyecto es $T_e = \max(EF_F) = \mathbf{20\text{ semanas}}$.
    \end{itemize}
  \item \textbf{Backward Pass (Tiempos Tardíos):} Recorremos el proyecto de atrás hacia adelante para determinar el término más tardío ($LF$) y el inicio más tardío ($LS$), fijando el tiempo final igual a la duración esperada ($LF_F = T_e = 20$).
    \begin{itemize}
      \item Para actividades finales: $LF = T_e$.
      \item Para cualquier otra actividad $i$: $LF_i = \min_{k} (LS_k)$ de todos sus sucesores inmediatos $k$.
      \item Inicio tardío: $LS_i = LF_i - t_e(i)$.
    \end{itemize}
  \item \textbf{Holguras ($H_i$) y Ruta Crítica:} La holgura total de una actividad representa el margen de retraso disponible sin aplazar el proyecto: $H_i = LF_i - EF_i = LS_i - ES_i$. Aquellas actividades con holgura nula ($H_i = 0$) son las actividades críticas.
\end{itemize}

El cálculo ordenado se resume en la siguiente tabla de tiempos de calendario:

\begin{center}
\small
\begin{tabular}{@{}lcccccccc@{}}
  \toprule
  \textbf{Actividad} & \textbf{Predecesores} & \textbf{$t_e$ (sem)} & \textbf{ES} & \textbf{EF} & \textbf{LS} & \textbf{LF} & \textbf{Holgura ($H$)} & \textbf{¿Crítica?} \\
  \midrule
  A & --   & 3 & 0  & 3  & 0  & 3  & 0 & Sí \\
  B & A    & 4 & 3  & 7  & 6  & 10 & 3 & No \\
  C & A    & 7 & 3  & 10 & 3  & 10 & 0 & Sí \\
  D & B, C & 6 & 10 & 16 & 10 & 16 & 0 & Sí \\
  E & B    & 5 & 7  & 12 & 11 & 16 & 4 & No \\
  F & D, E & 4 & 16 & 20 & 16 & 20 & 0 & Sí \\
  \bottomrule
\end{tabular}
\end{center}

\textbf{Ruta Crítica:} \textbf{A - C - D - F}. Duración esperada = \textbf{20 semanas}.

\textbf{3) Probabilidades e Intervalos de Confianza:}
\begin{itemize}
  \item Varianza de la duración total ($\sigma_T^2$):
\end{itemize}
    $$\sigma_T^2 = \sigma_A^2 + \sigma_C^2 + \sigma_D^2 + \sigma_F^2 = \frac{1}{9} + \frac{16}{9} + 1,0 + \frac{4}{9} = \frac{30}{9} \approx 3,333 \text{ semanas}^2$$
    $$\sigma_T = \sqrt{3,333} \approx 1,826 \text{ semanas}$$
\begin{itemize}
  \item Probabilidad de completar en menos de 22 semanas:
\end{itemize}
    $$Z = \frac{22 - 20}{1,826} = 1,095 \approx 1,10 \implies P(Z \le 1,10) = 0,8643 \implies \mathbf{86,43\%}$$
\begin{itemize}
  \item Intervalo de confianza del 95\%:
\end{itemize}
    $$[20 - 1,96(1,826); \; 20 + 1,96(1,826)] \implies [16,42; \; 23,58] \text{ semanas}$$

\textbf{4) Crashing (Aceleración de Proyecto):}
Queremos bajar de 20 a 18 semanas (reducción de 2 semanas) al costo mínimo acortando actividades de la ruta crítica:
\begin{itemize}
  \item \textbf{A:} Reducción máx 1 sem, costo marginal = \$3.000/semana.
  \item \textbf{C:} Reducción máx 2 sem, costo marginal = \$1.500/semana.
  \item \textbf{D:} Reducción máx 2 sem, costo marginal = \$2.500/semana.
  \item \textbf{F:} Reducción máx 1 sem, costo marginal = \$1.000/semana.
\end{itemize}

Estrategia de acortamiento óptimo:
1.  Acortamos \textbf{F} por 1 semana (costo = \$1.000). Duración baja a 19 semanas.
2.  Acortamos \textbf{C} por 1 semana (costo = \$1.500). Duración baja a 18 semanas.
\begin{itemize}
  \item Costo de crashing total = \$2.500.
  \item \textit{Decisión:} Conviene acortar \textbf{F por 1 semana y C por 1 semana}, obteniendo una ganancia neta de \$8.000 - \$2.500 = \$5.500.
\end{itemize}
\end{ejercicio}

\subsection{Last Planner (planificación operativa)}

\begin{definicion}{El Último Planificador (Ballard, 2000)}
El gran problema operativo: ``nunca resulta exactamente lo planificado'' (variabilidad + errores). El \textbf{Last Planner} cambia la visión \textit{push} por \textit{pull} e incorpora pasos intermedios: análisis de restricciones, compromisos entre las partes, mantener trabajos ``por hacer'' (\textit{lookahead}) y retroalimentación. Lógica: \textbf{Puede $\to$ Debe $\to$ Hace}. Lleva la planificación a la realidad: \textit{Planificación $\to$ Restricciones $\to$ Variabilidad}.
\end{definicion}

\subsection{V/F frecuentes de Proyectos}

\begin{center}
\small
\begin{tabular}{@{}clp{10.5cm}@{}}
  \toprule
  $\#$ & R & Afirmación y corrección \\
  \midrule
  1 & \textbf{V} & ``Las actividades de la ruta crítica tienen holgura cero.'' \\
  2 & \textbf{F} & ``La varianza del proyecto es la suma de las varianzas de todas las actividades.'' --- Falso: solo las de la \textbf{ruta crítica}. \\
  3 & \textbf{F} & ``Acortar cualquier actividad reduce la duración del proyecto.'' --- Falso: solo acortar actividades de la \textbf{ruta crítica} (y hasta que cambie la RC). \\
  4 & \textbf{V} & ``En PERT el tiempo esperado de una actividad no coincide con el tiempo más probable $m$.'' \\
  5 & \textbf{F} & ``Una pseudo-actividad consume recursos y tiene duración positiva.'' --- Falso: duración 0 y no consume recursos. \\
  \bottomrule
\end{tabular}
\end{center}

% ════════════════════════════════════════════════════════════════════════════




\end{document}
